\subsection{首批图表与表格说明：从协议层到正文层的落图规则}\label{subsec:typed-address-biaxial-completion-window6-figure-table-guidelines}
上一节已经把 window-\(6\) 的最小实例簇压成六层审计协议，并把“哪些量必须先被刚性核验、哪些量可以进入解释层”的逻辑顺序固定下来。
图表设计若要与这一协议相容，也必须服从同样的门槛链：先给固定窗闭式锚点，再给 boundary holonomy 证书，再给预算裂分，随后才允许进入表观随机与粗糙可见量。
否则读者会先看到“现象图”，却无法区分哪些量是闭式不变量，哪些量只是某种读出协议下的二次产物。

window-\(6\) 之所以适合作为首批图表的核心，并不在于它已经给出最终物理图景，而在于它已同时固定
\[
|X_6|=21,
\qquad
2{:}8,\ 3{:}4,\ 4{:}9,
\qquad
g_6=\frac{39}{64}\log 2+\frac{13}{64}\log 3,
\qquad
\kappa_6=\frac{11}{8}\log 2+\frac{3}{16}\log 3,
\]
\[
\dim_{\FF_2}\bigl(G_6^{\mathrm{bin}}\bigr)^{\mathrm{ab}}=21,
\qquad
X_6^{\mathrm{bdry}}=\{\texttt{100001},\texttt{100101},\texttt{101001}\},
\qquad
C_{\partial}\cong (\ZZ/2\ZZ)^3,
\]
以及精确回放阈值与 anomaly 完备预算之间的严格分裂。
因此，首批图表的目标不是“丰富视觉材料”，而是把这些刚性对象按协议顺序从实现层送到正文层。

\subsubsection{图组与表组的总顺序}\label{subsubsec:typed-address-biaxial-completion-window6-figure-table-order}
首批图表的排序不应由视觉效果决定，而应由审计协议的依赖关系决定。

\leanverified{paper\_typed\_address\_biaxial\_completion\_window6\_figure\_table\_order}
\begin{proposition}[首批图表与表格的落图顺序]\label{prop:typed-address-biaxial-completion-window6-figure-table-order}
围绕 window-\(6\) 的首批图表，宜按下列顺序进入正文：
\begin{enumerate}
  \item 先给可见层与纤维谱锚点图，并与刚性不变量总表并列；
  \item 再给 boundary 二层覆盖与最小 holonomy 图，并与 boundary/holonomy 审计表配对；
  \item 然后给回放预算与 anomaly 预算裂分图，并与预算等级—脚本映射表配对；
  \item 在前三层锁定后，才进入固定窗 collision、粗糙可见量采样以及块碰撞/mixed-collision 判别图；
  \item 最后再给跨尺度 boundary uplift 的协议门图，用以限制外推。
\end{enumerate}
\end{proposition}

\begin{proof}
这一顺序只是把上一节的六层协议改写成落图规则。
固定窗闭式、boundary holonomy 与预算裂分分别为统计层、粗糙读出层与跨尺度外推提供前提；若把后三类图提前，则其所依赖的刚性对象尚未锁定，图像便失去证据资格而退化为未经审计的现象展示。证毕。
\end{proof}

\subsubsection{可见层与纤维谱锚点图}\label{subsubsec:typed-address-biaxial-completion-window6-figure-visible-spectrum-anchor}
首幅图应承担的任务，不是展示“一个有限状态空间”，而是把整个最小实例的可见层锚死。
其题名宜固定为“window-\(6\) 的可见层、纤维多重度谱与规范体积锚点”。

该图宜分成两栏。
左栏展示
\[
X_6=X_6^{\mathrm{cyc}}\sqcup X_6^{\mathrm{bdry}},
\qquad
|X_6|=21,
\qquad
|X_6^{\mathrm{cyc}}|=18,
\qquad
|X_6^{\mathrm{bdry}}|=3,
\]
并把 boundary 三词
\[
X_6^{\mathrm{bdry}}
=
\{\texttt{100001},\texttt{100101},\texttt{101001}\}
\]
直接标出。
右栏展示 bin-fold 的固定窗纤维直方图
\[
2{:}8,
\qquad
3{:}4,
\qquad
4{:}9,
\]
并在图内同步标出
\[
g_6\approx 0.645542,
\qquad
\kappa_6\approx 1.159067,
\qquad
\dim_{\FF_2}\bigl(G_6^{\mathrm{bin}}\bigr)^{\mathrm{ab}}=21.
\]
该图只承担枚举与闭式核验功能，不承担拟合、插值或延拓功能。

正文在引用此图时，宜把它与可见层非平凡性绑在同一句中：window-\(6\) 的可见层并不是单一壳，而是一个由 \(18\) 个 cyclic 稳定词与 \(3\) 个 boundary 稳定词组成的刚性分裂；同时纤维谱 \(2{:}8,3{:}4,4{:}9\) 与 \((g_6,\kappa_6,\dim_{\FF_2}(G_6^{\mathrm{bin}})^{\mathrm{ab}})\) 三元锚点同步闭合，因而该窗口已经足以承载局部可逆、异常签名与高阶碰撞的第一批非平凡结构。

\subsubsection{boundary 二层覆盖与最小 holonomy 图}\label{subsubsec:typed-address-biaxial-completion-window6-figure-boundary-holonomy}
第二幅图应把 boundary 结构从“附加位”提升为最小 holonomy 层。
其题名宜固定为“boundary 二层覆盖、canonical sheet parity 与最小 holonomy”。

该图宜分成三栏。
左栏列出三条 boundary 词及其 dyadic 二层预像对，并在每条纤维旁标出固定差值
\[
34=F_9.
\]
中栏把各条 boundary 纤维对应的中心 involution 画成 parity 立方体，明确展示
\[
C_{\partial}\cong (\ZZ/2\ZZ)^3.
\]
右栏则单独标出几何可实现子群
\[
P_6^{\mathrm{geo}}=\langle(1,1,1)\rangle\cong \ZZ/2\ZZ,
\]
以及不可平坦化余商
\[
C_{\partial}/P_6^{\mathrm{geo}}\cong (\ZZ/2\ZZ)^2,
\qquad
r_{\mathrm{tor}}(Z_6^{\mathrm{bd}})=3.
\]
图注不应把这一层说成“也许存在某种相位”，而应明确指出：在最小窗口上，局域几何至多吸收同步对角方向，而仍有两位 residual freedom 不能被局域平坦化。

正文引用此图时，宜只服务于一个命题：boundary residual 不是局部噪声，而是局部可读—全局不可平的第一个有限窗证书。

\subsubsection{回放预算与 anomaly 预算裂分图}\label{subsubsec:typed-address-biaxial-completion-window6-figure-budget-split}
第三幅图的任务，是把账本预算从单一“位数”拆成任务等级。
其题名宜固定为“精确回放预算、boundary 超选预算与 anomaly 完备预算的有限窗裂分”。

该图应把两类对象放在同一坐标系中。
第一类对象是 fixed-window 预言机容量
\[
C_6(B)=8\min(2,2^B)+4\min(3,2^B)+9\min(4,2^B),
\]
它给出
\[
C_6(0)=21,
\qquad
C_6(1)=42,
\qquad
C_6(B)=64\quad(B\ge 2),
\]
从而最小精确回放阈值满足
\[
B_6^{\mathrm{replay}}=2.
\]
第二类对象不是容量曲线，而是最小完备异常预算
\[
B_6^{\mathrm{anom}}=21.
\]
若版面允许，还应在同图中显式标出
\[
B_6^{\mathrm{bdry}}=3,
\qquad
B_6^{\mathrm{anom}}-B_6^{\mathrm{replay}}=19.
\]
这条差额不是“需要更多比特”的模糊说法，而是“局部精确可逆”与“保留全部异常签名”之间的严格有限窗溢价。

因此，图注应把账本预算写成任务分层
\[
L\in\{0,1,2,21,+\infty\},
\]
而不应压成单一位数。
正文引用此图时，也应只服务于一个命题：可逆性与异常可见性不是同一语义任务，任何把二者混并的叙事，都会在 window-\(6\) 上立即失真。

\subsubsection{固定窗 collision 与均匀基线偏差图}\label{subsubsec:typed-address-biaxial-completion-window6-figure-fixed-window-collision}
进入统计层后，第一张图仍应从 fixed-window 闭式开始，而不是直接给长窗回归。
其题名宜固定为“window-\(6\) 的固定窗碰撞锚点与均匀基线偏差”。

该图宜分成两栏。
左栏展示精确碰撞常数
\[
S_2^{\mathrm{bin}}(6)=8\cdot 2^2+4\cdot 3^2+9\cdot 4^2=212,
\qquad
\mathrm{Col}_2^{\mathrm{bin}}(6)=\frac{212}{64^2}=\frac{53}{1024}.
\]
右栏展示相对于均匀基线 \(u_6\) 的二阶偏差
\[
\|p_6-u_6\|_2^2=\frac{89}{21504},
\qquad
\chi^2(p_6\|u_6)=\frac{89}{1024}.
\]
若版面允许，可在图下方附一个小条形图，直接拆分 \(d=2,3,4\) 三层对二点碰撞的贡献，以便把“高阶统计为何向顶端纤维冻结”直接落成可视结构。

正文引用此图时，应把它作为“表观随机”叙述的入口条件：window-\(6\) 的输出并非均匀打散后的轻微扰动，而是一个由大纤维显著抬升的可见分布；正因为这种抬升是闭式可算的，后面才有资格把“像随机”严格区分为由纤维谱导致的随机外观，而不是真噪声。

\subsubsection{粗糙可见量的最小采样图}\label{subsubsec:typed-address-biaxial-completion-window6-figure-rough-visibility}
粗糙可见量的第一张图必须避免把“看起来很粗糙的曲线”误当成证据。
其题名宜固定为“最小粗糙可见量 \(Y_\omega^{(\le J)}\) 的三层采样图”。

该图的采样协议应与定义 \ref{def:typed-address-biaxial-completion-window6-audit-rough-sampling} 完全一致，固定
\[
Y_\omega^{(\le J)}(t)
=
\sum_{j=0}^{J} a^j
\cos\!\Bigl(
2^j t+\phi_j^{\mathrm{vis}}+\phi_j^{\mathrm{bdry}}+\phi_j^{\mathrm{ledger}}
\Bigr),
\qquad 0<a<1,
\]
并同时绘出三条轨道：
\begin{enumerate}
  \item 只保留 \(\phi_j^{\mathrm{vis}}\) 的参考轨道；
  \item 加入 \(\phi_j^{\mathrm{bdry}}\) 的 boundary-holonomy 轨道；
  \item 再加入 \(\phi_j^{\mathrm{ledger}}\) 的全机制轨道。
\end{enumerate}
每条轨道都应附带局部放大窗，以检查“放大后是否继续出现新振荡”。
若局部放大窗只在第三条轨道上显著增粗，则说明真正驱动粗糙性的不是可见稳定层，而是账本/holonomy 的逐尺度回写；反之，若三条曲线几乎重合，则当前相位注入规则尚未把 residual 真正送回可见层。

因此，这张图没有资格脱离前四层单独出现。
它的输入坐标必须被明确声明为来自已经通过审计的 \(X_6\)、\(X_6^{\mathrm{bdry}}\)、\(C_{\partial}\)、预算等级 \(L\) 以及动态 prime-register 前缀。

\subsubsection{块碰撞缺口与 mixed-collision 图}\label{subsubsec:typed-address-biaxial-completion-window6-figure-block-collision}
与粗糙可见量配对的统计图，不应再回到一体分布，而应直接进入块层与跨机制比较。
其题名宜固定为“块碰撞缺口 \(\Delta_q^{(\ell)}\) 与归一化 mixed-collision 相似度 \(\Xi_q^{(0,1)}\)”。

该图宜采用双图布局。
左图画
\[
\Delta_q^{(\ell)}
:=
\log \mathrm{Col}_q^{(\ell)}
-\ell\log \mathrm{Col}_q^{(1)}
\]
随 \(q\) 或 \(\ell\) 的变化。
若该量始终接近 \(0\)，则块层几乎可因子化；若它在高 \(q\) 区间显著偏离，则说明确定性粗糙的时间结构已经进入可见块统计。

右图画
\[
\Xi_q^{(0,1)}
:=
\frac{\mathrm{Col}_q^{(0,1)}}
{\bigl(\mathrm{Col}_q^{(0,0)}\mathrm{Col}_q^{(1,1)}\bigr)^{1/2}}
\]
在三组比较中的走势：
deterministic rough 对 IID baseline，deterministic rough 对 visible-only baseline，以及 deterministic rough 对 boundary-only baseline。
这张图的目标不是判断“哪一种机制更乱”，而是回答：哪一种相位来源真正改变了尾部几何，哪一种只改变了一体边缘。

由于 \(\mathrm{Col}_q^{(0,1)}\) 在 finite-state online 机制下仍有限态可编译、可满足线性递推并由 Perron 根控制，这一图应被视为主统计图，而不是附加实验图。

\subsubsection{跨尺度 boundary uplift 的协议门图}\label{subsubsec:typed-address-biaxial-completion-window6-figure-crossscale-gate}
跨尺度图必须放在图组末端，因为它的功能不是展示新现象，而是阻止误外推。
其题名宜固定为“boundary sheet parity 的跨尺度不稳定性与协议门 \(\Gamma_m\)”。

图中应并列展示
\[
\Delta_6^{\mathrm{bdry}}=\{0,34\},
\qquad
\Delta_7^{\mathrm{bdry}}=\{0,55,89\},
\qquad
\Delta_8^{\mathrm{bdry}}=\{0,89,144\}.
\]
这样读者可以一眼看到：window-\(6\) 的二层覆盖并不会自动延拓到 \(m=7,8\)；同一 dyadic uplift 口径下，覆盖层数已经从二层变成三层。
因此，若要继续谈 sheet parity 的跨尺度保真，就必须额外定义并核验协议门
\[
\Gamma_m:\Delta_m^{\mathrm{bdry}}\to \widehat{\Delta}_m^{\mathrm{bdry}}.
\]
在没有 \(\Gamma_m\) 的情况下，任何跨尺度 parity 叙述都只应被视为未经许可的外推。

这张图的图注不应把重点放在“奇怪现象”，而应明确写出其限制性功能：window-\(6\) 的 boundary 二层覆盖是一个最小锚点事实，而不是自动跨尺度成立的规律。

\subsubsection{刚性不变量总表}\label{subsubsec:typed-address-biaxial-completion-window6-table-rigid-invariants}
与首幅图配对的总表，应把实现层最先需要复现的对象压成同一页。
其题名宜固定为“window-\(6\) 的刚性不变量总表”。

该表至少应设置四列：对象、精确值、语义层、审计入口。
其中最少应包含以下行：
\[
|X_6|=21,
\qquad
|X_6^{\mathrm{cyc}}|=18,
\qquad
|X_6^{\mathrm{bdry}}|=3,
\]
\[
\text{fiber histogram}=2{:}8,\ 3{:}4,\ 4{:}9,
\]
\[
\log |G_6^{\mathrm{bin}}|
=
39\log 2+13\log 3
\approx 41.314700,
\qquad
g_6\approx 0.645542,
\qquad
\kappa_6\approx 1.159067,
\]
\[
\dim_{\FF_2}\bigl(G_6^{\mathrm{bin}}\bigr)^{\mathrm{ab}}=21,
\qquad
\mathrm{Col}_2^{\mathrm{bin}}(6)=\frac{53}{1024}.
\]
该表中的 \(g_6\)、\(\kappa_6\)、异常秩与 fixed-window 纤维直方图共同构成整组图表的停机线：任何后续图若与此表不一致，都应被直接判定为实现层错误。

\subsubsection{boundary 与 holonomy 审计表}\label{subsubsec:typed-address-biaxial-completion-window6-table-boundary-holonomy}
第二张表应把最小曲率候选的离散证书集中到一页上。
其题名宜固定为“boundary residual 与 holonomy 审计表”。

该表宜设置五列：对象、精确值、最小几何解释、是否局域可吸收、备注。
其中至少应包含
\[
X_6^{\mathrm{bdry}}=\{\texttt{100001},\texttt{100101},\texttt{101001}\},
\qquad
\delta_{\mathrm{bdry}}=34=F_9,
\]
\[
C_{\partial}\cong (\ZZ/2\ZZ)^3,
\qquad
P_6^{\mathrm{geo}}=\langle(1,1,1)\rangle\cong \ZZ/2\ZZ,
\]
\[
C_{\partial}/P_6^{\mathrm{geo}}\cong (\ZZ/2\ZZ)^2,
\qquad
r_{\mathrm{tor}}(Z_6^{\mathrm{bd}})=3.
\]
该表的用途不是重复正文，而是为后文任何关于 phase residual、几何 uplift 或局域吸收的论证提供一张可直接调用的离散证书页。

\subsubsection{预算等级与脚本映射表}\label{subsubsec:typed-address-biaxial-completion-window6-table-budget-script-map}
第三张表应把数学量、任务语义与实现入口钉在一起。
其题名宜固定为“预算等级、任务语义与脚本映射表”。

该表应至少包含如下等级：
\[
L=0
\quad
(\text{同图表单值读出}),
\qquad
L=1
\quad
(\text{单条 boundary 纤维选层}),
\]
\[
L=2
\quad
(\text{window-}6\text{ 精确回放}),
\qquad
L=21
\quad
(\text{异常签名完备}),
\qquad
L=+\infty
\quad
(\text{当前协议下无忠实外置}).
\]
对应的审计入口至少应列出
\[
\texttt{scripts/exp\_fold\_bin\_gauge\_m6\_invariants.py},
\qquad
\texttt{scripts/exp\_window6\_c6\_orbit\_patisalam\_seed.py},
\]
\[
\texttt{scripts/exp\_fold\_collision\_moment\_recursions.py},
\qquad
\texttt{scripts/exp\_fold\_collision\_pressure\_multifractal\_q60.py},
\]
并为后续粗糙可见性、块碰撞、boundary holonomy 与 prime-ledger 审计保留
\[
\texttt{scripts/exp\_window6\_rough\_visibility\_audit.py},
\qquad
\texttt{scripts/exp\_window6\_block\_collision\_audit.py},
\]
\[
\texttt{scripts/exp\_window6\_boundary\_holonomy\_audit.py},
\qquad
\texttt{scripts/exp\_window6\_prime\_ledger\_prefix\_audit.py}
\]
四个专题入口。
这张表的价值在于：它把“数学量—语义层—实现层”三者固定在同一张映射表中，从而避免后续讨论只剩抽象符号而没有相应的审计入口。

\subsubsection{正文中的引用规则}\label{subsubsec:typed-address-biaxial-completion-window6-figure-table-citation-rule}
图表一旦进入正文，其引用方式也必须与协议顺序保持一致。

\begin{proposition}[正文中的图表引用应采用 claim--evidence 结构]\label{prop:typed-address-biaxial-completion-window6-figure-table-citation-rule}
\leanverified{paper\_typed\_address\_biaxial\_completion\_window6\_figure\_table\_citation\_rule}
在结果段中，围绕 window-\(6\) 的首批图表，正文宜按下列顺序组织：
\begin{enumerate}
  \item 先用可见层与纤维谱锚点图及刚性不变量总表，证明最小窗口已经是非平凡多壳结构；
  \item 再用 boundary 二层覆盖图与 boundary/holonomy 审计表，证明 boundary residual 不是局部噪声，而是不可全部几何吸收的 holonomy；
  \item 继而用预算裂分图与预算—脚本映射表，证明可逆预算与异常预算严格分裂；
  \item 在前三层锁定后，再进入 fixed-window collision、粗糙可见量采样与块碰撞/mixed-collision 图，讨论表观随机与粗糙几何；
  \item 最后用跨尺度协议门图收口，明确任何跨尺度 parity 论断都需要额外协议门。
\end{enumerate}
\end{proposition}

\begin{proof}
这一引用规则只是命题 \ref{prop:typed-address-biaxial-completion-window6-figure-table-order} 在正文层的改写。
前三层图表分别固定可见层、residual 层与预算层；后两类图若脱离这些锚点被单独引用，则统计或几何信号将失去证据资格。证毕。
\end{proof}

换言之，首批图表的职责不是“让结果更直观”，而是把解释层的自由度锁定在实现层已经通过的硬对象之后。
即使读者不同意更高层的物理解释，他也必须先面对这些有限窗对象本身：纤维直方图、中心子群、预算裂分、collision 常数与跨尺度失败证书。
解释可以争论，锚点不能漂移。

\subsubsection{小结}\label{subsubsec:typed-address-biaxial-completion-window6-figure-table-summary}
因此，围绕 window-\(6\) 的首批图表不应被看作“附加展示”，而应被看作从协议层进入正文层的第一组证据载体：
可见层图与总表固定锚点，
boundary 图与 holonomy 表固定 residual，
预算图与脚本映射表固定任务等级，
统计图与粗糙采样图才在此前提下承担表观随机与粗糙可见性的说明功能，
而跨尺度协议门图则负责限制所有未经许可的外推。
下一步最自然的工作，便是把这组图表进一步压成逐段可复用的结果段落模板，使图、表与闭式证书能够直接回写到正文叙述中。
